\documentclass{beamer}
\usepackage{moreverb} 
\usepackage{listings}
\usepackage{mflogo}
% imprimir
% \documentclass[handout]{beamer} 
% \usepackage{pgfpages}
% \pgfpagesuselayout{4 on 1}[a4paper,landscape,border shrink=5mm]

\mode<presentation> {
  \usetheme{Warsaw}
  \setbeamercovered{opaque}
}

\usebackgroundtemplate{\includegraphics[width=\paperwidth]{format/libresoft-bg.png}}
\usepackage[spanish]{babel}
\usepackage[utf8]{inputenc}
\usepackage{graphics}
\usepackage{amssymb} % Simbolos matematicos

%% Metadatos del PDF.
\hypersetup{  
  pdftitle={Virtualización con software libre},
  pdfauthor={Miguel Vidal},
  pdfcreator={GSyC/Libresoft},
  pdfproducer=PDFLaTeX,
  pdfsubject={Master on Free Software},
}
%%

\defbeamertemplate*{footline}{shadow theme}
{%
  \leavevmode%
  \hbox{\begin{beamercolorbox}[wd=.5\paperwidth,ht=2.5ex,dp=1.125ex,leftskip=.1cm plus1fil,rightskip=.2cm]{author in head/foot}%
\includegraphics[scale=0.40]{format/cc-by-80x15.png} \hspace{0.05cm}
	Miguel Vidal / Jose Castro
  \end{beamercolorbox}%
  \begin{beamercolorbox}[wd=.5\paperwidth,ht=2.5ex,dp=1.125ex,leftskip=.3cm,rightskip=.3cm plus1fil]{title in head/foot}%
    \usebeamerfont{title in head/foot}\insertshorttitle%
    \hspace{2cm} \usebeamerfont{author in head/foot}\insertframenumber\,/\,\inserttotalframenumber\hfill
  \end{beamercolorbox}}%
  \vskip0pt%
}


\begin{document}

\title{Virtualización con software libre}
\subtitle{System Integration}
% \institute{\texttt{http://gsyc.urjc.es/\~{}mvidal} \\ Twitter: \texttt{@mvidallopez}}
\author{Miguel Vidal, Jose Castro} 
\date{\footnotesize{Master on Free Software \\ April 13rd, 2012}}
\author{Miguel Vidal \hspace{1cm} Jose Castro \\
\hspace{0.5mm} {\tiny Twitter: @mvidallopez \hspace{1.1cm}Twitter: @jfcastroluis}
}


\frame{
\maketitle
\begin{center}
\includegraphics[width=6cm]{format/gsyc-urjc}
\end{center}
}

%% License slide
\begin{frame}
  \vspace{2cm}
  \begin{flushright}
    {\small \copyright{} 2010-2012 Miguel Vidal, Jose Castro} \\
%    \vspace{0.25cm}
    \medskip
    {\scriptsize This work is licensed under \\ a Creative Commons Attribution 3.0 License}
%    \vspace{0.10cm}
  \end{flushright}
  \begin{flushright}
    \href{http://creativecommons.org/licenses/by/3.0/es}{\includegraphics[width=2cm]{format/cc-by.png}} \\
    {\tiny \url{http://creativecommons.org/licenses/by/3.0}}
  \end{flushright}
\end{frame}%%

\usebackgroundtemplate{}

\AtBeginSection[]
{
\begin{frame}<beamer>
\begin{center}
{\huge \insertsection}
\end{center}
\end{frame}
}


\AtBeginSubsection[]
{
  \begin{frame}<beamer>{Índice}
    \tableofcontents[currentsection,currentsubsection]
  \end{frame}
}

%%


%%%%%%%%%%%%%%%%%%%%%%%%%%%%%%%%%%%%%%%%%%%%%%%%%%%%%%%%%%%%%%%%%%%%%%%
\section{Qué es la virtualización}
%%%%%%%%%%%%%%%%%%%%%%%%%%%%%%%%%%%%%%%%%%%%%%%%%%%%%%%%%%%%%%%%%%%%%%%

\begin{frame}
\frametitle{Qué es la virtualización}

% \begin{definition}
\begin{block}{Definición}
Es la combinación de hardware y software que permite a un \alert{recurso físico} (como un servidor, un sistema operativo o un dispositivo de almacenamiento) funcionar como \alert{múltiples recursos lógicos}.
% \end{definition}
\end{block}

\begin{itemize}
\item También se puede definir como la \alert{abstracción} de los recursos de una máquina.
\end{itemize}

\end{frame}

%%%%%%%%%%%%%%%%%%%%%%%%%%%%%%%%%%%%%%%%%%%%%%%%%%%%%%%%%%%%%%%%%%%%%%%

\subsection{Orígenes}

\begin{frame}
\frametitle{Virtualización: Orígenes remotos}

\begin{itemize}
\item Término acuñado en los años 60 (también conocido como ``pseudo máquina'').
\item El proyecto M44/44X (mediados de los 60): el objetivo era evaluar los conceptos de los sistemas de tiempo compartido.
\item En los 70 se desarrollaron varios sistemas de máquinas virtuales basadas en IBM.
\item Un componente llamado \alert{Virtual Machine Monitor (VMM)} corría directamente sobre el hardware ``real''.
\item Se podían crear varias máquinas virtuales y cada instancia ejecutaba su propio sistema operativo.
\end{itemize}

\end{frame}

%%%%%%%%%%%%%%%%%%%%%%%%%%%%%%%%%%%%%%%%%%%%%%%%%%%%%%%%%%%%%%%%%%%%%%%

\begin{frame}
\frametitle{Antecedentes modernos}

\begin{itemize}
\item \alert{1999}: VMware introduce el primer producto de virtualización x86, VMware Virtual Platform.
\item \alert{2005}: VMWare distribuye gratuitamente virtualización para escritorio.
\item \alert{Finales de 2005}: Intel introduce VT-x. \alert{2006}: AMD-V.
\item \alert{2006}: Parche de Xen para kernel Linux 2.6.18.
\item \alert{2007}: KVM aparece en la rama oficial del kernel Linux 2.6.20.
\item \alert{2007}: Sun publica LDOM para Sparc. 
\end{itemize}

\end{frame}

%%%%%%%%%%%%%%%%%%%%%%%%%%%%%%%%%%%%%%%%%%%%%%%%%%%%%%%%%%%%%%%%%%%%%%%

\subsection{Conceptos básicos}


%%%%%%%%%%%%%%%%%%%%%%%%%%%%%%%%%%%%%%%%%%%%%%%%%%%%%%%%%%%%%%%%%%%%%%%

\begin{frame}
\frametitle{Conceptos básicos}

\begin{itemize}
\item \alert{anfitrión} (\textit{host}): el SO que ejecuta el software de virtualización.
	\begin{itemize}
	\item El anfitrión controla el hardware real.
	\end{itemize}
\item \alert{invitado} o \alert{huésped} (\textit{guest}): el SO virtualizado.
	\begin{itemize}
	\item Puede haber varios invitados en un mismo anfitrión. 
	\item Los anfitriones no deben interferir entre ellos ni con el anfitrión.
	\end{itemize}
\end{itemize}

\end{frame}

%%%%%%%%%%%%%%%%%%%%%%%%%%%%%%%%%%%%%%%%%%%%%%%%%%%%%%%%%%%%%%%%%%%%%%%


\begin{frame}
\frametitle{Conceptos de virtualización}

\begin{itemize}
\item Al software de virtualización se le llama:
	\begin{itemize}
	\item  \alert{Hipervisor}.
	\item  \alert{Virtual Machine Manager} or \alert{VMM}.
	\end{itemize}
\item El VMM o Hipervisor corre como parte del sistema operativo del anfitrión (o es el anfitrión).
\item A una instancia del hardware virtualizado se la conoce como \alert{Máquina Virtual} o \alert{VM}.
\item Los sistemas operativos huéspedes corren dentro de una VM.
\end{itemize}

\end{frame}


%%%%%%%%%%%%%%%%%%%%%%%%%%%%%%%%%%%%%%%%%%%%%%%%%%%%%%%%%%%%%%%%%%%%%%%

\begin{frame}
\frametitle{Hipervisor (I)}

\begin{itemize}
\item Los \textit{Virtual Machine Monitors} (\textit{aka} \alert{hipervisores}) permiten que diferentes SOs, tareas y configuraciones de software coexistan en una misma máquina física.
\item Abstrae los recursos físicos de la máquina anfitriona para las distintas ``máquinas virtuales''.
\item Garantizan nivel de aislamiento entre los invitados.
\item También proporcionan una interfaz única para el hardware.
\end{itemize}

\end{frame}

%%%%%%%%%%%%%%%%%%%%%%%%%%%%%%%%%%%%%%%%%%%%%%%%%%%%%%%%%%%%%%%%%%%%%%%

\begin{frame}
\frametitle{Hipervisor (y II)}

Dos clases de hipervisores:

\begin{itemize}
\item \alert{Tipo 1} (o ``nativo'', ``bare-metal''): el hipervisor es una capa entre el hardware y el sistema operativo. 
\begin{itemize}
\item Al sistema operativo se le llama Dominio de Control  y corre sobre el hipervisor.
\item Los guests son ``Dominios Lógicos''.
\end{itemize}

\item \alert{Tipo 2} (o ``hosted''): el hipervisor es una capa de software que corre sobre el sistema operativo anfitrión.
\end{itemize}

\end{frame}


%%%%%%%%%%%%%%%%%%%%%%%%%%%%%%%%%%%%%%%%%%%%%%%%%%%%%%%%%%%%%%%%%%%%%%%

\begin{frame}
\frametitle{Extensiones de virtualización para x86}

\begin{itemize}
\item 2005: Intel y AMD añaden soporte hardware para la virtualización.
\item Intel Virtualization Technology (VT) \textit{codename} Vanderpool
\item AMD Virtualization (AMD-V) \textit{codename} Pacifica
\item Añaden una funcionalidad específica para permitir a los hipervisores un rendimiento mayor en virtualización completa.
\item La virtualización completa es más sencilla de implementar.
\end{itemize}

\end{frame}



%%%%%%%%%%%%%%%%%%%%%%%%%%%%%%%%%%%%%%%%%%%%%%%%%%%%%%%%%%%%%%%%%%%%%%%

\subsection{Razones}

\begin{frame}
\frametitle{Razones para virtualizar (1)}

\begin{itemize}
\pause
\item Permite el \alert{aislamiento} de aplicaciones y usuarios sobre la misma máquina para que no se interfieran entre sí.
\item \alert{Consolidación}: fusiona varios servidores o servicios en una misma máquina (email, web, dns, etc.).
\item Ejecutar diferentes sistemas operativos o software específico para un tipo de CPU.
\end{itemize}

\end{frame}

%%%%%%%%%%%%%%%%%%%%%%%%%%%%%%%%%%%%%%%%%%%%%%%%%%%%%%%%%%%%%%%%%%%%%%%

\begin{frame}
\frametitle{Razones para virtualizar (y 2)}

\begin{itemize}
\item \alert{Entornos de prueba} sencillos y seguros: depuración, desarrollo y aislamiento para estudiar virus, etc.
\item \alert{Fácil desarrollo} de software utilizando aplicaciones de software virtuales. Desde la perspectiva del negocio, es posible reducir los costes totales de propiedad (TCO).
\item \alert{Minimiza el consumo de energía} y la infraestructura de refrigeración en los \textit{datacenters}.
\end{itemize}

\end{frame}



%%%%%%%%%%%%%%%%%%%%%%%%%%%%%%%%%%%%%%%%%%%%%%%%%%%%%%%%%%%%%%%%%%%%%%%

\begin{frame}
\frametitle{Virtualización y Cloud Computing}

\begin{itemize}
\item El Cloud Computing (IaaS) \alert{no} es lo mismo que la gestión de la virtualización.
\item Pero es cierto que la mayoría de los entornos de Cloud Computing hacen uso de la virtualización.
\item La gestión del hardware está totalmente abstraída.
\item La capacidad de la infraestructura en Cloud Computing es elástica (tanto para crecer como para menguar).
\end{itemize}

\begin{center}
\alert{El cloud (IaaS) son recursos hardware convertidos en un \\ ``\textit{pool} de recursos''}
\end{center}


\end{frame}



%%%%%%%%%%%%%%%%%%%%%%%%%%%%%%%%%%%%%%%%%%%%%%%%%%%%%%%%%%%%%%%%%%%%%%%
\section{Tipos de virtualización}
%%%%%%%%%%%%%%%%%%%%%%%%%%%%%%%%%%%%%%%%%%%%%%%%%%%%%%%%%%%%%%%%%%%%%%%


\begin{frame}[<+->]
\frametitle{Tipos de virtualización}

El objetivo es crear la impresión de tener hardware separado en un único sistema físico. Hay \alert{4 maneras} de hacer esto:

\begin{enumerate}
\item \alert{Emulación}
\item \alert{Virtualización completa}
\item \alert{Paravirtualización} 
\item \alert{Virtualización a nivel de sistema operativo}
\end{enumerate}

\pause
¿Alguno más?

\end{frame}

%%%%%%%%%%%%%%%%%%%%%%%%%%%%%%%%%%%%%%%%%%%%%%%%%%%%%%%%%%%%%%%%%%%%%%%

\subsection{Emulación}

\begin{frame}
\frametitle{Emulación}

\begin{itemize}
\item La máquina virtual simula el hardware completo.
\item Huéspedes sin modificar de diferentes arquitecturas hardware corren dentro de una VM.
\item Utilizado para \alert{sandboxes} o para desarrollar software sin que el hardware esté disponible físicamente.
\item \alert{Ejemplos:} Qemu; basilisk II (un emulador de Mac (m68k) para x86).
\end{itemize}
\end{frame}

\begin{frame}
\frametitle{Emulación: arquitectura}
\begin{center}
\includegraphics[scale=0.35,clip=false]{figs/emulation.png}
\end{center}

\begin{itemize}
\item \alert{Ventajas:} Simular hardware que no está físicamente disponible.
\item \alert{Desventajas:} Bajo rendimiento y baja densidad (coste alto).
\end{itemize}


\end{frame}


%%%%%%%%%%%%%%%%%%%%%%%%%%%%%%%%%%%%%%%%%%%%%%%%%%%%%%%%%%%%%%%%%%%%%%%

\subsection{Completa}


\begin{frame}
\frametitle{Virtualización completa}

\begin{itemize}
\item Parecido a la emulación, anfitriones sin modificar.
\item Se diferencia de la emulación en que los sistemas operativos a virtualizar están diseñados para correr en la misma arquitectura que el anfitrión.
\item Combinado con hardware: CMT, IntelVT, AMD-V, estas CPUs controlan el acceso a instrucciones de virtualización.
\item \alert{Ejemplos:} KVM, VirtualBox, Xen + Intel VT. 
\end{itemize}
\end{frame}

%%%%%%%%%%%%%%%%%%%%%%%%%%%%%%%%%%%%%%%%%%%%%%%%%%%%%%%%%%%%%%%%%%%%%%%

\begin{frame}
\frametitle{Virtualización completa: arquitectura}
\begin{center}
\includegraphics[scale=0.35,clip=false]{figs/fullvirt.png}
\end{center}

\begin{itemize}
\item \alert{Ventajas:} Flexibilidad, ejecuta diferentes sistemas operativos de distintos fabricantes.
\item \alert{Desventajas:} No se pueden emular otras arquitecturas. Rendimiento (en CPUs estándar).
\end{itemize}

\end{frame}


%%%%%%%%%%%%%%%%%%%%%%%%%%%%%%%%%%%%%%%%%%%%%%%%%%%%%%%%%%%%%%%%%%%%%%%

\subsection{Paravirtualización}

\begin{frame}
\frametitle{Paravirtualización}

\begin{itemize}
\item El hipervisor ofrece una versión modificada del anfitrión.
\item La VM ofrecida tiene la misma arquitectura que el anfitrión.
\item \alert{Ejemplo:} Xen con CPUs estándar.
\end{itemize}
\end{frame}

\begin{frame}
\frametitle{Paravirtualización: arquitectura}
\begin{center}
\includegraphics[scale=0.35,clip=false]{figs/paravirt.png}
\end{center}

\begin{itemize}
\item \alert{Ventajas:} Rendimiento (ligero y rápido), escalabilidad y facilidad de gestión; fuerte aislamiento; permite virtualizar sin necesidad de CPUs especiales. 
\item \alert{Desventajas:} Para la misma arquitectura. Requiere modificar el SO del invitado.
\end{itemize}


\end{frame}



%%%%%%%%%%%%%%%%%%%%%%%%%%%%%%%%%%%%%%%%%%%%%%%%%%%%%%%%%%%%%%%%%%%%%%%

\subsection{Ligera}

\begin{frame}
\frametitle{Virtualización ligera}

\begin{itemize}
\item También conocida como \alert{virtualización a nivel de SO} o \textit{virtual hosts}.
\item La virtualización se hace instanciando la imagen del SO (sin hipervisor).
\item El SO está modificado para permitir múltiples procesos en diferentes espacios de usuario aislados unos de otros, cada uno con su configuración de red.
\item \alert{Ejemplos:} Jaulas de FreeBSD, zonas/contenedores de OpenSolaris, Linux VServers. 
\end{itemize}

\end{frame}


\begin{frame}
\frametitle{Virtualización ligera: arquitectura}
\begin{center}
\includegraphics[scale=0.50,clip=false]{figs/armitage2.png}
\end{center}

\begin{itemize}
\item \alert{Ventajas:} Rápida, capa ligera de virtualización. Rendimiento cercano al nativo. Densidad.
\item \alert{Desventajas:} Es difícil implementar el aislamiento fuerte. Tampoco se pueden (a priori) virtualizar distintos SO.
\end{itemize}

\end{frame}



%%%%%%%%%%%%%%%%%%%%%%%%%%%%%%%%%%%%%%%%%%%%%%%%%%%%%%%%%%%%%%%%%%%%%%%

\begin{frame}
\frametitle{Otros tipos de virtualización (1)}

\begin{itemize}
\item \alert{Virtualización de bibliotecas:} biblioteca \textbf{Wine} (subconjunto de la API de Win32 para poder ejecutar aplicaciones Windows)
\item \alert{Virtualización de aplicación:} entorno de ejecución virtual (con una API para la ejecución en diferentes plataformas). \textit{Ejemplo:} \textbf{Java Virtual Machine}.
\item \alert{Virtualización de escritorio:} se implementa el escritorio como servicio. \textit{Ejemplo:} \textbf{SunVDI}.
\end{itemize}

\end{frame}

%%%%%%%%%%%%%%%%%%%%%%%%%%%%%%%%%%%%%%%%%%%%%%%%%%%%%%%%%%%%%%%%%%%%%%%

\begin{frame}
\frametitle{Otros tipos de virtualización (y 2)}

\begin{itemize}
\item \alert{Virtualización de almacenamiento:} \textit{arrays} y \textit{pools} de discos. \textit{Ejemplos:} \textbf{LVMs} en Linux y \textbf{ZFS} en OpenSolaris. También los FS distribuidos (\textbf{OCFS2}, \textbf{GFS}).
\item \alert{Virtualización de red:} \textbf{Crossbow} (nativa de OpenSolaris) permite crear switches e interfaces de red virtuales.
\item Memoria virtual, particiones de bases de datos...
\end{itemize}

\end{frame}


\section{Tecnologías de virtualización}

%%%%%%%%%%%%%%%%%%%%%%%%%%%%%%%%%%%%%%%%%%%%%%%%%%%%%%%%%%%%%%%%%%%%%%%

\subsection{Xen}


\begin{frame}
\frametitle{Xen: Características (1)}

\begin{itemize}
\item Xen utiliza \alert{paravirtualización}.
\item El sistema operativo huésped tiene que estar modificado para usar el hipervisor.
\item También admite virtualización completa (con Intel VT o AMD-V): 
	\begin{itemize}
	\item huéspedes sin modificar 
	\item permite que sistemas operativos privativos (como Windows) puedan ser virtualizados.
	\end{itemize}
\end{itemize}

\end{frame}

%%%%%%%%%%%%%%%%%%%%%%%%%%%%%%%%%%%%%%%%%%%%%%%%%%%%%%%%%%%%%%%%%%%%%%%

\begin{frame}
\frametitle{Xen: Características (2)}

\begin{itemize}
\item Es capaz de hacer \alert{migración} de máquinas virtuales.
\item Xen (HV) se ejecuta en el anillo de protección 0 mientras que los dominios se ejecutan en el anillo 1 o anillo 3.

\begin{figure}
  \centering
\includegraphics[scale=0.40,clip=false]{figs/rings-xen.png}
  \caption{\footnotesize{Anillos: dominios de protección jerárquica (x86)}}
\end{figure}

\end{itemize}

\end{frame}


%%%%%%%%%%%%%%%%%%%%%%%%%%%%%%%%%%%%%%%%%%%%%%%%%%%%%%%%%%%%%%%%%%%%%%%

\begin{frame}
\frametitle{Xen: Prestaciones (1)}
\begin{itemize}

\item \alert{Independencia} entre los sistemas virtualizados. Se puede reiniciar y crear independientemente.
\item \alert{Uso mejorado del hardware}: balanceo de recursos. Una máquina virtual puede hacer uso de los recursos que no utilizan las otras máquinas virtuales.
\item \alert{Backup sencillo}. Sólo con copiar la máquina virtual se puede levantar en un nuevo servidor. Xen también permite la migración en caliente, siendo muy flexible y minimizando el tiempo de recuperación en caso de fallo.
\end{itemize}

\end{frame}

%%%%%%%%%%%%%%%%%%%%%%%%%%%%%%%%%%%%%%%%%%%%%%%%%%%%%%%%%%%%%%%%%%%%%%%

\begin{frame}
\frametitle{Xen: Prestaciones (2)}
\begin{itemize}

\item Se pueden \alert{modificar parámetros} como la RAM, el número de CPUs y el espacio en disco para cada necesidad específica de cada máquina virtual.
\item \alert{Entornos de prueba} y desarrollo: múltiples máquinas virtuales en un único servidor físico para probar y desarrollar.
\end{itemize}

\end{frame}


%%%%%%%%%%%%%%%%%%%%%%%%%%%%%%%%%%%%%%%%%%%%%%%%%%%%%%%%%%%%%%%%%%%%%%%

\begin{frame}
\frametitle{Xen: Cómo funciona}
\begin{center}
\includegraphics[width=9cm]{figs/XEN-schema.png}
\end{center}
\end{frame}

%%%%%%%%%%%%%%%%%%%%%%%%%%%%%%%%%%%%%%%%%%%%%%%%%%%%%%%%%%%%%%%%%%%%%%%

\begin{frame}
\frametitle{¿Xen integrado en Linux?}

\vspace{-0.25cm}
\begin{center}
\includegraphics[width=5.5cm]{figs/xen-kernel-support.png} \\
\end{center}

\begin{flushright}
\includegraphics[width=5.5cm]{figs/xen-kernel-support3.png} \\
\end{flushright}

\begin{center}
\includegraphics[width=5cm]{figs/xen-kernel-support2.png} 
\end{center}

\end{frame}


%%%%%%%%%%%%%%%%%%%%%%%%%%%%%%%%%%%%%%%%%%%%%%%%%%%%%%%%%%%%%%%%%%%%%%%

\begin{frame}
\frametitle{Xen: Linux 3.0}
\begin{itemize}

\item En Linux 3.0 los \alert{drivers paravirtualizados de Xen} se integran en el kernel oficialmente.
\item Ahora pueden usarse Dom-0 (host) y DomU sin modificar/hackear el kernel.
\item El hipervisor Xen sigue siendo un proyecto desarrollado aparte de Linux.
\item KVM sigue siendo el único hipervisor integrado en el kernel Linux.
\end{itemize}

\end{frame}

%%%%%%%%%%%%%%%%%%%%%%%%%%%%%%%%%%%%%%%%%%%%%%%%%%%%%%%%%%%%%%%%%%%%%%%

\begin{frame}
\frametitle{Xen en kernel Linux 2.x}
\begin{center}
\includegraphics[width=9cm]{figs/xen-old-kernels.png}
\end{center}
\end{frame}

%%%%%%%%%%%%%%%%%%%%%%%%%%%%%%%%%%%%%%%%%%%%%%%%%%%%%%%%%%%%%%%%%%%%%%%

\begin{frame}
\frametitle{Xen en kernel Linux 3.x}
\begin{center}
\includegraphics[width=9cm]{figs/xen-new-kernel.png}
\end{center}
\end{frame}


%%%%%%%%%%%%%%%%%%%%%%%%%%%%%%%%%%%%%%%%%%%%%%%%%%%%%%%%%%%%%%%%%%%%%%%

\subsection{KVM}


\begin{frame}
\frametitle{Linux: Kernel-based Virtual Machine (1)}

\begin{itemize}
\item KVM es un infraestructura de virtualización específica del kernel Linux.
\item Virtualización nativa (completa) utilizando los micros Intel VT y AMD-V.
\item La primera versión fue incluida en Linux 2.6.20 (febrero 2007).
\item Ha sido portada por Joyent a illumos (fork libre de OpenSolaris).
\end{itemize}

\end{frame}

%%%%%%%%%%%%%%%%%%%%%%%%%%%%%%%%%%%%%%%%%%%%%%%%%%%%%%%%%%%%%%%%%%%%%%%

\begin{frame}
\frametitle{Linux: Kernel-based Virtual Machine (2)}

\begin{itemize}
\item Muchos SOs funcionan virtualizados con KVM: Linux, BSD, Solaris, Windows, Haiku, ReactOS y  MacOS X.
\item No requiere emulación: un programa en espacio de usuario (Qemu) utiliza la interfaz \texttt{/dev/kvm} para configurar el espacio de direcciones de las VMs hospedadas.
\end{itemize}

\end{frame}

%%%%%%%%%%%%%%%%%%%%%%%%%%%%%%%%%%%%%%%%%%%%%%%%%%%%%%%%%%%%%%%%%%%%%%%

\begin{frame}
\frametitle{KVM: arquitectura}

\begin{center}
\includegraphics[scale=0.40,clip=false]{figs/kvm_architecture.png}
\end{center}

Dos componentes:
\begin{enumerate}
\item Un driver del kernel Linux (\texttt{/dev/kvm}).
\item Un componente en espacio de usuario para emular el hardware del PC (Qemu).
\end{enumerate}
\end{frame}

%%%%%%%%%%%%%%%%%%%%%%%%%%%%%%%%%%%%%%%%%%%%%%%%%%%%%%%%%%%%%%%%%%%%%%%

\begin{frame}
\frametitle{KVM: Virtualización completa en Linux}
\begin{center}
\includegraphics[scale=0.2,clip=false]{figs/kvm.png}
\end{center}

\end{frame}

%%%%%%%%%%%%%%%%%%%%%%%%%%%%%%%%%%%%%%%%%%%%%%%%%%%%%%%%%%%%%%%%%%%%%%%

\begin{frame}
\frametitle{KVM: SmartOS (illumos)}
\vspace{-0.25cm}
\begin{figure}
\includegraphics[width=10cm]{figs/KVM-xp-booted.png}
\caption{KVM en SmartOS, virtualizando Windows XP.}
\end{figure}

\end{frame}


%%%%%%%%%%%%%%%%%%%%%%%%%%%%%%%%%%%%%%%%%%%%%%%%%%%%%%%%%%%%%%%%%%%%%%%

\subsection{Jails}


\begin{frame}
\frametitle{FreeBSD: Jails}

\begin{center}
\includegraphics[scale=0.50,clip=false]{figs/armitage3.png}
\end{center}

\begin{itemize}
\item Instancias (réplicas) del kernel FreeBSD aisladas, cada una con su IP, cuentas de usuario, etc.
\item Virtualización ligera (un solo kernel).
\end{itemize}

\end{frame}



%%%%%%%%%%%%%%%%%%%%%%%%%%%%%%%%%%%%%%%%%%%%%%%%%%%%%%%%%%%%%%%%%%%%%%%

\subsection{OpenSolaris}

\begin{frame}
\frametitle{Proyectos xVM OpenSolaris}

\begin{itemize}
\item \alert{VirtualBox}: VMs para todo. Diferentes plataformas.
\item \alert{xVM server}: Xen en Opensolaris 2009.06.
\item Dominios lógicos (\alert{LDOMs}): Hipervisor tipo 1, ``completo'', para plataformas SPARC.
\item \alert{Zonas}: virtualización ligera para OpenSolaris.
\item Otros: Sun XVM Ops Center, Sun xVM VDI.
\end{itemize}

\end{frame}

%%%%%%%%%%%%%%%%%%%%%%%%%%%%%%%%%%%%%%%%%%%%%%%%%%%%%%%%%%%%%%%%%%%%%%%

\begin{frame}
\frametitle{Contenedores y zonas}

\begin{itemize}
\item Tipo de virtualización que ejecuta varias instancias del mismo sistema operativo (y el mismo kernel).
\item Zona Global y Zonas Locales: no se puede acceder al sistema de ficheros global (un \texttt{chroot} mejorado).
\item Sin hipervisor ni hardware específico para virtualización.
\item Despliegue y administración muy simple: \texttt{zonecfg} y \texttt{zoneadm}.
\item x86: con soporte para ``branding''.
\item Desde 2005 forma parte de Solaris 10 (y OpenSolaris).
\item Parecido a las jaulas de FreeBSD, OpenVZ (contenedores Linux) o Linux-VServer.
\end{itemize}

\end{frame}


%%%%%%%%%%%%%%%%%%%%%%%%%%%%%%%%%%%%%%%%%%%%%%%%%%%%%%%%%%%%%%%%%%%%%%%

\begin{frame}
\frametitle{Virtualización con LDOMs}

\begin{itemize}
\item ``Virtualización completa'', basada en un hipervisor tipo 1.
\item Requiere CPUs especiales (Chip Multithreading = CMT).
\item Cada dominio es una máquina virtual completa con un conjunto configurable de recursos.
\item Los sistemas operativos que corren dentro del Dominio Lógico pueden iniciarse, pararse y reiniciarse de forma independiente.
\item OS base: Solaris 10 / Opensolaris 2009.06
\item OS de máquina virtual: Solaris 10, Opensolaris 2009.06, Sparc Linux y otros que soporten estan arquitectura.
\item Pero por ahora, tan solo SunOS está soportado.
\end{itemize}

\end{frame}


%%%%%%%%%%%%%%%%%%%%%%%%%%%%%%%%%%%%%%%%%%%%%%%%%%%%%%%%%%%%%%%%%%%%%%%

\begin{frame}
\frametitle{Chip Multithreading (CMT)}

\begin{itemize}
\item UltraSparc T1 / T2.
\item \alert{Multithread}: Un \textit{thread} es similar a una CPU.
\item Ejemplo: T1 tiene 8 cores con 4 threads/core.
\item SSL directamente soportado por el hardware (1 MAU/core).
\item LDOMs pueden asignar cores a las VMs.
\item El hipervisor ejecuta en el firmware del servidor.
\item Free/Open Hardware: \url{http://www.opensparc.net}
\item Sun Fire T / Enterprise T / Blade T servers.
\end{itemize}

\begin{center}
\includegraphics[width=6cm,clip=false]{figs/sparcT1.png}
\end{center}

\end{frame}

%%%%%%%%%%%%%%%%%%%%%%%%%%%%%%%%%%%%%%%%%%%%%%%%%%%%%%%%%%%%%%%%%%%%%%%


\begin{frame}
\frametitle{Referencias}

\begin{itemize}

\item Amit Singh ``An Introduction to Virtualization'' (2004): \\
\url{http://www.kernelthread.com/publications/virtualization/}
\item Comparison of platform virtual machines: \url{http://en.wikipedia.org/wiki/Comparison_of_platform_virtual_machines}
\item Jeanne Matthews \textit{et al.} \textit{Running Xen}. Prentice Hall, 2008 (Chapter 1: ``Background and Virtualization Basics'')

\end{itemize}

\end{frame}

%%%%%%%%%%%%%%%%%%%%%%%%%%%%%%%%%%%%%%%%%%%%%%%%%%%%%%%%%%%%%%%%%%%%%%%


\end{document}



